% ============================================================
% ch1_introduction.tex — 第一章 绪论
% ============================================================
\chapter{绪论}

\section{研究背景}

长上下文推理把大语言模型（Large Language Model, LLM）的部署压力集中到解码阶段。LLaMA~3~\cite{grattafiori2024llama3}、Qwen2.5~\cite{qwen2025qwen25} 等模型已经能够处理更长输入，但在自回归生成中，每个新 token 都要访问此前累积的缓存。对本文而言，部署压力的直接落点是：当上下文长度和服务并发同步上升时，Decode 阶段需要反复读取的 KV Cache 字节数会怎样限制可部署性。

Transformer~\cite{vaswani2017attention} 自回归解码的计算路径决定了这一压力的来源。第 $t$ 个生成步只新增当前位置的 Query、Key 和 Value，但注意力仍需访问前 $S$ 个历史位置的 Key/Value。KV Cache 避免了重复执行历史 K/V 投影，却把缓存容量与访存带宽变成每个 decode step 都要面对的约束。也就是说，缓存不是只在 prefill 后静态占用显存；它还会在后续 token 生成中持续参与注意力计算。

从存储结构上看，KV Cache 的显存占用与模型层数、每层 KV 头数、头维度、序列长度以及存储精度近似成正比，可写为
\begin{equation}
  M_{\mathrm{KV}} = 2 \times L \times H_{kv} \times d_k \times S \times b
  \label{eq:kv_memory}
\end{equation}
其中 $L$、$H_{kv}$ 与 $d_k$ 由模型结构给定，$S$ 对应上下文长度，$b$ 对应缓存元素的存储字节数，系数 2 来自 Key 与 Value 两部分。本文使用式~\eqref{eq:kv_memory} 固定后续实验变量：模型结构决定基础缓存形状，序列长度与 batch size 决定部署压力，量化位宽直接改变 $b$，而预算分配方法则决定哪些层或角色保留更高精度。在长上下文和批量推理场景下，这一开销会随上下文长度与并发批次同步放大。vLLM~\cite{kwon2023vllm} 等推理框架通过分页注意力（PagedAttention）改善显存碎片和缓存管理效率，但这类方法并不直接减少 KV Cache 的绝对存储量。

因此，KV Cache 量化的系统价值首先来自它直接作用于 $b$。INT8 可以提供相对保守的缓存压缩锚点，INT4 则进一步压低存储字节数，也更容易暴露低比特扰动在注意力计算中的后果。本文后续不把“位宽更低”等同于“路径更好”，而是把 Decode 阶段缓存字节数、注意力行为稳定性和任务读数放在同一套实验口径下比较。

\section{问题定义与研究动机}

在上述背景下，KV Cache 量化的价值首先来自直接的存储压缩。INT8 可以作为较保守的保真基线，更低位宽则能进一步释放显存空间，支持更长上下文或更高并发。但缓存压缩不能只按 bit 数下降来理解：位宽降低以后，Key 与 Value 的扰动会进入注意力计算，并可能改变检索、摘要和语言建模等任务行为。本文后续以 INT4 作为低比特压力点，正是因为这一设置会更清楚地暴露 Key/Value 角色差异和模型间敏感度差异；后文的 Qwen 与 LLaMA 对照会把这种差异落到 \texttt{K4V8}/\texttt{K4V4} 等布局上，此处只把它作为问题线索，并保留架构调制的解释空间。

当前 KV Cache 量化的校准目标仍常落在数值重建误差上，典型做法包括 MSE 损失和基于分位数的裁剪。这类做法通常把“张量越接近原值，模型行为越接近全精度路径”作为工作前提。进入注意力计算后，这个前提并不总稳固：softmax 的非线性可能把较小的 Key 扰动转化为注意力概率的明显变化，而某些数值上更大的误差又可能在归一化中被部分吸收。因此，张量误差大小与下游功能偏移之间难以建立稳定的一一对应关系；不同层、不同头和不同上下文位置上的分布差异，也会进一步削弱静态阈值的解释力。

KV Cache 量化的困难因此不只是继续调节超参数或压低重建误差。模型在推理时感知到的是量化误差经过注意力计算后的功能偏移：注意力分布是否被重排，聚合输出是否偏离，任务行为是否随之改变。特别是在长上下文和低比特场景下，同一量级的数值扰动落在不同层、不同头或不同角色上，可能产生完全不同的后果。单纯追求全局数值误差最小，并不能保证模型在推理时保持稳定和可用。

这使本文把观察对象从单纯的张量近似扩展到推理行为本身：KV Cache 量化在保持重建精度之外，还需要尽量维持注意力分布、聚合输出和任务表现的稳定。量化后注意力分布是否发生重排、这种重排是否继续影响输出与任务结果，是更贴近推理可用性的观察依据。后续关于校准目标、量化路径和预算分配的讨论，也以这些行为读数作为共同依据。

为避免“行为”停留在抽象表述，本文将注意力行为操作化为三类可观测量：其一是量化前后注意力分布的偏移，用于刻画检索和排序关系是否被扰动；其二是加权聚合输出的变化，用于刻画 Value 信息进入残差流后的表示偏移；其三是语言建模、检索和长文档任务上的表现变化，用于检验上述扰动是否传导到任务层面。进一步地，本文所说的行为敏感度画像，是指由离线校准与诊断读数聚合得到的逐层相对敏感度，用于回答哪些层或角色应被优先保护；策略适用区间则指在特定模型家族、规模、任务和预算条件下，不同预算策略进入高性能区域的条件范围。

基于这一问题意识，下一节只做必要的相关研究定位：已有工作分别回答低比特表示、缓存压缩格式和系统执行效率问题，而本文要接上的缝隙是量化缓存进入注意力计算后如何保持推理行为。

\section{国内外研究现状概述}

围绕长上下文推理与 KV Cache 压缩，现有研究主要回答三个相邻但不同的问题。模型量化研究提供低比特表示、校准与误差控制的基础方法；KV Cache 专属量化直接处理缓存存储压力；高效推理系统则从注意力计算和缓存管理角度降低部署开销。三条路线的交集说明，长上下文部署由压缩精度、缓存组织和系统执行共同决定。

在模型量化方向，已有研究围绕权重与激活的低比特表示、训练后量化以及校准策略形成了较为成熟的方法体系。以 SmoothQuant~\cite{xiao2023smoothquant}、GPTQ~\cite{frantar2023gptq} 和 AWQ~\cite{lin2024awq} 为代表的工作，分别从量化难度转移、近似误差补偿和激活感知保护等角度推进了大模型的压缩与部署。这些方法主要面向模型参数与激活，但其中关于量化表示、误差控制和校准流程的基本思想，也为后续 KV Cache 压缩提供了重要的方法学参考。

在 KV Cache 专属量化方向，随着长上下文推理对缓存显存的压力持续上升，研究开始直接面向缓存本身设计压缩路径。KIVI~\cite{liu2024kivi} 强调 K/V 统计差异与异构量化轴，KVQuant~\cite{hooper2024kvquant} 关注极低比特表示中的码本和误差补偿，ZipCache~\cite{he2024zipcache} 则从重要性驱动的缓存压缩切入。这些工作留下的共同问题是：当缓存扰动通过 softmax 和加权聚合进入推理行为时，校准与格式选择应如何被审计。

除量化压缩外，高效推理方向也在同步推进。FlashAttention~\cite{dao2022flashattention,dao2024flashattention2} 通过更高效的分块注意力计算降低中间访存开销；分页注意力~\cite{kwon2023vllm} 则通过缓存分页改善显存碎片和缓存利用率。这类工作与 KV Cache 量化并非替代关系，而是从不同层面共同服务于长上下文推理的可部署性。

现有研究已经从模型量化、KV Cache 专属压缩与高效推理等角度推进了长上下文部署，但这些路线解决的层面并不相同。本文的切入点是 KV Cache 量化中的行为保持问题：量化后的缓存不仅要占用更少显存，还要在注意力计算中保持可用的推理行为。下面的研究问题从这一观察出发，而不是预设某一种量化格式或预算策略会在所有模型上占优。

\section{核心假设与研究问题}

KV Cache 量化的核心难点不只是继续提高压缩率，而是如何在压缩后保持模型的实际推理行为。本文把注意力行为作为分析起点，关注量化后注意力分布、聚合输出和任务表现是否仍与全精度路径保持一致。

\noindent\textbf{核心假设：} 在 KV Cache 量化中，注意力行为比张量重建误差更接近模型实际功能损伤。本文关心量化后注意力分布是否保持稳定、加权聚合输出是否偏离，以及这些变化是否影响任务行为。基于这一假设，校准目标和预算分配可以共享同一个观察对象，而不必被拆成彼此孤立的局部问题。

\noindent\textbf{研究问题 1：注意力行为能否作为 KV Cache 量化的统一对象？} KV Cache 量化需要保持的，到底只是张量层面的重建精度，还是还包括注意力分布、聚合输出和任务行为的稳定性？如果后者能够成立，它应当同时支撑基准校准、低比特恢复和预算分配三类讨论，而不是只解释某一个局部设置。

\noindent\textbf{研究问题 2：行为引导校准能否形成稳定的 INT8 基准路径？} 在 INT8 这一较保守、便于审计的位宽下，以注意力行为为导向的校准能否建立一条相对 FP16 任务行为偏差可控的基准路径？这一路径也为后续低比特探索和系统落地提供相对稳定的参照。

\noindent\textbf{研究问题 3：行为敏感度画像能否支撑跨模型预算分配？} 校准链路产生的行为敏感度画像，能否在不同模型家族、不同规模和不同任务条件下读出可区分的策略适用区间？本文关注的是这些条件如何改变预算策略的有效范围，而不是给出一种脱离条件的单点配置。

\noindent\textbf{研究问题 4：固定预算选择是否需要画像驱动的自动建议？} 手工固定 $k$ 的预算 sweep 在跨模型迁移中是否会暴露不稳定性？同源敏感度画像能否被转写为 \texttt{AutoK} 预算建议，从而减少手工选择带来的偶然性？本文把 \texttt{AutoK} 放在 allocation 的自动化扩展位置；它的经验作用由第四章的跨模型读数限定。

围绕上述核心假设与四个问题，下一节将进一步概述本文的研究内容、技术路线与主要贡献。

\section{研究内容、技术路线与主要贡献}
\label{sec:ch1-contribution}

围绕前述核心假设与四个问题，本文将 KV Cache 量化组织为一条以注意力行为为主线的研究路径。校准层回答如何选择量化参数，低比特路径回答如何在更激进的压缩下保持可用，分配层回答有限预算应优先保护哪些层或角色，预算建议层回答固定 $k$ 是否需要画像驱动的自动建议。

本文从注意力行为视角重新组织 KV Cache 量化问题，用一个更贴近模型实际功能稳定性的对象连接不同层级的量化决策；在此基础上，围绕不同位宽形成可执行的方法实例，并通过跨模型比较提炼策略适用区间。本文据此追问：校准、低比特实例和预算分配能否被同一组行为读数解释，而不把贡献限定在某个局部设置的压缩收益上。

第二章用于定位已有工作和本文边界，第三章给出行为代理、路径实例和分配机制，第四章围绕基准保真、低比特恢复、跨模型结构读数与 \texttt{AutoK} 预算建议给出实验证据，第五章收束结论与适用边界。

\noindent\textbf{以注意力行为保持重新界定 KV Cache 量化对象。} 校准目标从单一张量重建误差扩展到注意力分布、聚合输出和任务行为这组可观测量。校准层与预算分配层因此读取同一组行为证据，用于判断量化扰动是否进入推理功能，而不是各自优化互不相连的局部指标。

\noindent\textbf{建立从 INT8 基准路径到低比特恢复的可复核实例链。}
INT8 基准路径用于检验行为引导校准在保守位宽下的基础保真能力；
进入低比特场景后，\mbox{INT4-RoleAlign} 用 Key 与 Value 的功能角色差异组织失稳诊断与路径实例；
\texttt{AutoK} 则把同一行为敏感度画像转写为预算建议。
Triton~\cite{tillet2019triton} kernel、离线校准产物和在线推理管线构成系统支撑层，
使这条路径能在当前实验协议下执行，并提供固定复核接口。

\noindent\textbf{识别跨模型预算策略的适用区间。} 跨模型实验不把某一种策略写成所有模型家族、规模和任务条件下的统一最优，而是用来刻画同量级 INT4 预算带内的策略适用区间。这样的读法为后文分析 heuristic 强基线、关键层偏离、\texttt{AutoK} 模型级候选、早层保护与高性能簇共存等现象提供统一口径；这些现象的成立范围由第四章实验读数限定。

上述内容与四个高层问题对应：先确认注意力行为是否足以组织 KV Cache 量化，再建立 INT8 保真锚点；低比特恢复用于检验同一对象在更激进位宽下是否仍能指导路径实例，随后跨模型预算分配与 \texttt{AutoK} 建议共同回答画像读数能否进入预算决策。

\section{论文结构安排}
\label{sec:ch1-structure}

全文按上述问题展开。第一章界定 Decode 阶段 KV Cache 压力、行为保持假设和四个研究问题。第二章说明 Transformer、KV Cache、量化与高效推理中哪些已有结论可直接复用，哪些问题仍未覆盖缓存扰动后的注意力行为。第三章给出行为代理、INT8 基准路径、\mbox{INT4-RoleAlign} 和预算分配接口。第四章给出基准保真、低比特恢复、跨模型策略和 \texttt{AutoK} 预算建议的实验证据，并把部署效率读数作为系统边界而非单独的加速宣称。第五章收束贡献、适用范围和未覆盖的任务与部署条件。
